%!TEX root = ../thesis.tex
%************************* Final Chapter ***************

\chapter{Conclusions}
\label{chap:Conc}

\graphicspath{{Chap_Conclusion/Figs/}}
	
	In this thesis we have presented and examined a number of methods related to finding Early Warning Signals for tipping points in dynamical systems. The methods we have chosen for study are mostly extensions of the ``fingerprinting'' technique of \cite{held2004} in which the value of a certain `indicator' is tracked over time and a change in this value is supposed to be a warning of a tipping point. Thus, the indicator value, plotted as a function of time, is the early warning \emph{signal}. In general, it is not the actual value that matters but the change in the value.
	
	In chapter~\ref{chap:1D} we have studied the uses of the ACF1 indicator which is common in the literature, the DFA indicator which was developed by \cite{livina2007}, and the PS indicator which is a result of this work \citep{prettyman2018,prettyman2019}. All of these indicators have been found, in certain systems containing tipping points, to provide useful early warning signals, although in these cases the data from which the EWS is obtained is a one-dimension time series. A pressing problem in the field of early warning signals is, therefore, to find methods which can be applied to higher-dimensional systems such as those presented by \cite{williamson2015} and \cite{williamson2016} and the analysis of spatial correlations by \cite{bathiany2013p2} or \cite{kefi2014early}. In the climate science in particular the dynamical systems under investigation are often extremely high-dimensional. The methods presented in chapter~\ref{chap:2D} are intended as an important step in this direction, although we restrict our study to two-dimensional systems. The example in chapter~\ref{chap:cyclones}, that of a tropical cyclone, is likewise studied in the context of one- or two-dimensional early warning signal techniques, with the one-dimensional time series of sea-level pressure data being studied either alone or in conjunction with a second variable: wind speed. 
	
	In the remainder of this chapter we conclude this thesis by providing comments on the work we have done and extending it, hypothetically, into the future by considering lines of research which may take this thesis as a starting point. We do not attempt to summarise the entire thesis here but concentrate on the most significant results from each of the different paths we have travelled along our journey. 
	
	
	
\section{Comments on the PS indicator}
\label{sec:conc:PScomments}

The power spectrum scaling exponent seems an obvious candidate for an early warning indicator since it is closely related to the lag-1 autocorrelation and the DFA exponent which are already used as indicators. In chapter~\ref{chap:1D} we have introduced the idea of using the periodogram for the estimation of the PS exponent, first binning the fast Fourier transform periodogram to give an even distribution of points on the logarithmic scale. 

We base our understanding of tipping point detection and prediction on the idea of \emph{critical slowing down} (see section~\ref{subsec:1D:Exponents_as_EWS:CriticalSD}, page~\pageref{subsec:1D:Exponents_as_EWS:CriticalSD}), which is assumed to precede a tipping point, and which we have modelled as an AR(1) process. This assumption is used in section~\ref{subsec:1D:Exponents_as_EWS:determining_frequency_range} (page~\pageref{subsec:1D:Exponents_as_EWS:determining_frequency_range}) to determine a suitable range of frequencies in which to estimate the exponent. 

As with the development of the DFA exponent for use as an EWS indicator, the question arises that the PS exponent method may not be suitable when the power spectrum of a particular time series does not exhibit genuine power-law scaling. This objection is dealt with in section~\ref{subsec:1D:Exponents_as_EWS:nopowerlawsacaling} where we investigate the analytic relationship between the PS exponent and the parameter $\mu$ of the AR(1) process for which no real power-law scaling exists for $0<\mu<1$. We find that by estimating the PS exponent $\beta$ anyway, finding a linear best fit to the non-linear power spectrum `crossover' function, the result is a consistent, monotonically increasing value of $\beta$ which can be used as a proxy for increasing critical slowing down, similar to the DFA exponent and the lag-1 autocorrelation. In addition, we find the same thing when we assume the AR(1) parameter $\mu$ is non constant over the time series used to estimate the PS exponent (see section~\ref{subsec:1D:Exponents_as_EWS:PSnonstationary}, page~\pageref{subsec:1D:Exponents_as_EWS:PSnonstationary}), in this case the value of $\beta$ we estimate numerically is consistent with what we predict analytically by integrating the power spectrum of the AR(1) process over the range of $\mu$ values. 

Although this analysis confirms that the PS indicator may, in theory, provide an EWS for tipping points which are characterised by critical slowing down, it does not comment on the suitability of the indicator for tipping points which are \emph{not} so characterised, nor in cases when the critical slowing down cannot be modelled as an AR(1) process (the latter case has not been considered in this thesis). Experimentally, we find that the PS indicator does not provide an EWS for noise-induced tipping, consistent with established results for the ACF1 and DFA indicators. Transitional tipping points may or may not be characterised in this way. A future study may be designed to investigate these transitional tipping points further, and also to apply the PS indicator to examples of rate-induced tipping, which has not been studied here. 

Finally, we find that the estimation of the PS exponent requires a relatively long time series, in comparison to the accurate estimation of the lag-1 autocorrelation, for use as an EWS indicator. Where we consider examples of dynamical systems with genuine bifurcations, particularly the super-critical pitchfork bifurcation which is studied in more depth in section~\ref{subsec:1D:indicatorComparison:PitchforkModel}, we use a window size of only $100$ points for the PS indicator, which is not expected to produce a noticeable EWS, but a slight increasing trend can be often be observed nonetheless in the mean of a large enough ensemble. 

We can conclude that the use of the power spectrum scaling exponent as an EWS indicator is, in most cases, inferior to the well-established use of the simple lag-1 autocorrelation. There are, however, some aspects of this new method which may provide benefits and which could be exploited in future work. In particular, the PS exponent is insensitive to certain trends (although we note that the DFA exponent has already been exploited for this reason) and periodicities, particularly when the frequency of the periodicity lies outside the range of frequencies used for the PS exponent estimation. 

The application of the PS indicator to a two-variable system in chapter~\ref{chap:2D}, and to the tropical cyclone problem in chapter~\ref{chap:cyclones}, did not give much more insight into the PS indicator than the one-dimensional `toy model' examples in chapter~\ref{chap:1D}: once again we find that where the ACF1 indicator provides an EWS, so does the PS indicator, although not as strong nor as clear. When the ACF1 indicator fails to provide an EWS, so does the PS indicator. The exception to this is the application to sea-level pressure in the vicinity of tropical cyclones. In this case we discover that the PS indicator rises slightly before the tipping event, while the ACF1 indicator decreases. This inconsistency is possibly due to the 12-hour periodicity in the data due to tidal fluctuations, which is destroyed by the sudden large drop in pressure at the event time. When this periodicity was removed, using either a sine-wave subtraction or a deseasonalising approach, all EWS indicators became more variable and no clear signal was visible in any — it is likely that by subtracting a periodic function from the entire time series we are actually introducing a periodicity into that part of the time series which is not obviously fluctuating originally. In these situations, where attempting to remove periodicities or trends artificially might introduce biases, it may be useful to further explore the potential of the PS indicator which is insensitive to such artefacts. 


\section{Comments on the use of EOFs for dimension reduction}

In chapter~\ref{chap:2D} we have explored the idea of reducing the dimension of time series data using empirical orthogonal functions, thus allowing the use of established one-dimensional EWS techniques when higher-dimensional data is given. This idea is not original to this thesis and we take for granted that using EOFs to reduce dimensionality will give a better result than arbitrarily choosing to analyse, say, only $x$ coordinates. However, the question arises as to whether the EOF projection is optimal for use with EWS techniques. When attempting to find an EWS in a time series, one is interested in looking for \emph{an increase in variance} or \emph{an increase in autocorrelation}, etc.. But the EOF projection finds the time series with the largest \emph{variance}, not the largest \emph{increase in variance} or the largest \emph{increase in autocorrelation}. 

We have taken a simple, general, discrete-time, two-dimensional system as an example and attempted to discover whether the EOF method is optimal, or even appropriate, for use with EWS techniques. What we have discovered is that the EOF projection (which has maximal variance) is the same as, or very close to, the projection which has maximal \emph{increase in variance}, unless there is very large variance in a stationary part of the system. This result is, we believe, not trivial and serves to justify the use of EOFs in a range of EWS applications. 

An alternative EOF-like method for dimension reduction was also proposed, in which the lag-1 autocorrelation is maximised rather than the variance. We found that there was not a significant benefit to using this approach, even when the EWS indicator used on the resulting one-dimensional time series was the lag-1 autocorrelation itself. This result further justifies the use of the standard EOF method in EWS applications. 

We can conclude, therefore, that the EOF method is well-suited to reducing dimension prior to applying EWS techniques and the previously (to our knowledge) unexamined hypothesis that the maximal variance projection will be useful in an application which seeks to find an \emph{increase in variance}, is true in the cases we have considered. In future work it may be instructive to consider systems in dimensions higher than two and to consider other examples of systems which are generalisable to a wider variety of systems experiencing tipping points. 


\section{Comments on other two-dimensional methods}

In this thesis we have been particularly concerned with the method presented by \cite{williamson2015} (see section~\ref{sec:2D:Williamson}, page~\pageref{sec:2D:Williamson}) which estimates numerically the eigenvalues of the system Jacobian matrix in order to obtain a precursor signal of a bifurcation (for example, the negative-real eigenvalues moving towards a positive-real value). Because of the assumptions made on the dynamical system we view this method as similar to using lag-1 autocorrelation as an EWS based on the assumption of the presence of critical slowing down. We have used this method, therefore, as an example against which we compare the performance of the PS indicator in chapter~\ref{chap:2D}, in a natural extension of the work in chapter~\ref{chap:1D} where the PS and ACF1 indicators were compared directly in one-dimensional time series. 

Our results accurately recreated those of \cite{williamson2015} and we found that the EWS is as good as that produced by the ACF1 indicator coupled with the EOF method for dimension reduction, if one knows which eigenvalue to study and has some prior knowledge of how the EWS will look in the case of each specific system. 

In addition, we investigated a number of techniques for detecting tipping points in multiple time series over a 2D field. We presented a simple technique whereby an EWS indicator is evaluated in every time series over the field simultaneously to build up a heat map of a measure of increase in that indicator over a window of time. We find that the resulting picture of the propagation of a bifurcation point across the field is very clear when using the ACF1 indicator and accurately represents the state of the system at each point in time. 

Many other, similar methods exist in the literature (see \cite{bathiany2013p2}, \cite{kefi2014early}) but these were not studied further here, partly because the geophysical example chosen as a test-bed in chapter~\ref{chap:2D}, that of the approach of a tropical cyclone, did not present a good opportunity for studying these methods due to limited data. A future study which aims to test the PS indicator in other systems could, if the data existed in the correct form, use some of these other methods for a comparison. 





\section{Comments on applications to tropical cyclones}

The study of sea-level pressure data obtained from ground-based weather stations on the trajectory of a tropical cyclone is problematic due to sparse data sources and short time series. We were able to identify fourteen weather stations meeting the criteria of being in the path of a large tropical cyclone and having uninterrupted data at the time of the cyclone\footnote{Weather stations are often damaged during unusually powerful storms.}. In each of these fourteen cases we obtained hourly data, giving us only around $100$ relevant data points (often fewer) since a cyclone will often have progressed from first detection (from a given coastal location) to making landfall within a period of $100$ hours. This is just about the lower limit at which the PS and DFA indicators can be applied to a time series. When attempting to detect an effect which may only be detectable in a period of maybe $48$ hours, but the indicator is calculated in a window of $100$ hours, there will necessarily be a watering-down of the effect on the indicators. For this reason, we have not expected to discover any useful early warning signals, and it surprising that any change in the DFA and PS indicators is observed at all. It is not clear from the plots made of the single-value indicators applied to the sea-level pressure series, but when assessing the sensitivity of the indicators to window size we have presented contour plots, or heat maps, of the indicator value at different times before the cyclone event using different window sizes. In these plots we are able to see that the PS and DFA indicators do show a rising trend around $24$ hours before the tipping point when using larger window sizes (>80 points). This is maybe more noticeable in comparison to the same methods applied using shorter window sizes. It is somewhat surprising that there is any noticeable trend at all, no matter how slight, since the analysis of the PS exponent of the AR(1) process (see section~\ref{subsec:1D:Exponents_as_EWS:PSsensitivity}, page~\pageref{subsec:1D:Exponents_as_EWS:PSsensitivity}) does not suggest that the estimation of the PS exponent from just $100$ points would be at all consistent. It will be necessary, in future work, to apply the PS exponent to a real-life system with much more detailed datasets available, preferably a system which has previously been analysed using other EWS methods \citep{lenton2012}. 

The particular problem of approaching tropical cyclones was chosen as an example of a system with an obvious tipping point, in the sense of a sudden qualitative change, but for which the \emph{type of tipping} was unknown. It is unlikely that this example is representative of a genuine bifurcation and it is more likely to be a forced transition. Indeed, the model developed in section~\ref{sec:TC:Model} models the tipping event as a forced transition. Given this situation it was not known in advance whether the tipping point indicator methods would, in theory, provide an EWS. Nevertheless, we have shown that the PS exponent and DFA exponent are very slightly sensitive to the change in sea-level pressure in the presence of a tropical cyclone. A future study should be designed to establish in which situations a forced transition might be expected to exhibit critical slowing down or other effects to which established EWS methods are sensitive.

Another issue with this particular choice of system is the 12-hourly tidal oscillations in the sea-level pressure data which may affect the ACF1 and DFA indicators. When these oscillations were removed, either by deseasonalising or by subtracting a sine wave, the signals from the indicators became even less clear than was the case when using the raw, oscillating time series. Although this periodicity in the data might provide a useful test of the potential benefits of the PS indicator which we have shown to be insensitive to trends and periodicities, the necessarily short time series and small ensemble sizes involved when studying this particular system make the PS exponent a poor choice for an EWS indicator and outweigh the benefits. It would be interesting in further studies of the PS indicator to use additional examples of dynamical systems with periodic components, where better time series data is available than in this cyclone example, in order to establish whether the PS indicator could out-perform the simple ACF1 indicator in these situations. 

In addition, this choice of geophysical system did not provide a particularly good opportunity for testing EWS techniques which require multiple time series datasets spread over a 2D field (see section~\ref{sec:TC:field_data}, page~\pageref{sec:TC:field_data}). Few locations were identified where multiple weather stations with good data were available in the same geographic region when a tropical cyclone was also present. A coastal area of the Gulf of Mexico was selected, encompassing the coasts of Florida, Louisiana and Texas, where nine large  hurricanes\footnote{Tropical cyclones occurring in the North Atlantic ocean are known as hurricanes.} made landfall in the time period of the available data. This area contained $65$ weather stations with available data, although due to the variable quality of the data only between $21$ and $48$ stations were used in each case depending on the hurricane (see column ``\# stations'' in table~\ref{tab:TC:9hurricanes}, page~\pageref{tab:TC:9hurricanes}). In a future study the same methods could be applied to a system for which many hundreds of time series datasets are available at regular, gridded locations over a 2D field — possibly model data or satellite reanalysis data. One interesting subject which was explored preliminarily in this thesis is the dynamics of the \emph{greening of the Sahara}. We have looked into the work of, for example, \cite{bathiany2013p1} which uses model data with high spacial density to detect spacial correlations between geographic areas at discrete points in time. The same model with a higher time resolution is being developed \citep{bathiany2013p2}, which would allow the application of our higher-dimensional EWS indicator methods and therefore allow us to compare these methods to the spacial-correlation methods used previously. The data from this model is unavailable at the time of writing but this will make an interesting further study along the same lines as this thesis, possibly giving a better example than the tropical cyclone problem used here. 

In summary we can conclude that the tropical cyclone problem has some interesting aspects, for example it is a tipping point of unknown (probably non-bifurcational) mechanism, thus allowing to test our EWS techniques `blind', and the time series contain periodic oscillations which potentially allow the assessment of certain benefits of the PS indicator. However, the data we have used in our analysis does not allow for sufficiently long time series nor sufficiently large ensembles to properly compare our EWS techniques, nor does it allow for very in-depth study of the higher-dimensional techniques introduced in chapter~\ref{chap:2D}.

\section{Suggestions for future work}

One of the most pressing problems in the study of tipping points is that of predicting the time until the tipping occurs. Any method for predicting tipping points in this way would have to be informed by the nature of the dynamical system in question, and the type of tipping that is expected. It would be necessary, in advance, to classify a range of different varieties of tipping points in a range of different varieties of dynamical systems so that the precursors of the tipping points would be well-known. Our study of the various toy models used as examples throughout this thesis, in particular the super-critical pitchfork bifurcation and the general, discrete-time two-dimensional system in section~\ref{sec:2D:EOF-justification}, goes some way towards this. However, we find that, for continuous-time systems, changing the time-step for integration even slightly can significantly change the values of the various EWS indicators. 

In all of our examples we do not consider the actual values of the EWS indicators, nor a quantitative measure of their change over time (which will also be affected by the length of the sliding window), but only `whether it is increasing or not'. We are a long way from being able to say ``For a system that looks like \emph{this}, if the indicator is increasing like \emph{this}, then the tipping point will occur in $x$ time steps.''  And it seems this must be a long-term research goal.

However, we have identified, throughout this thesis and this concluding chapter, a number of suggestions for future work which may build upon the work of this thesis and strengthen the field of tipping point research. 

\begin{enumerate}
	\item {\bf Rate of change of parameters.} In chapter~\ref{chap:1D} we have explained our reasons for the choice of the parameters in each of the toy models we study using the EWS techniques (see section~\ref{subsec:1D:indicatorComparison:choice_of_parameters}, page~\pageref{subsec:1D:indicatorComparison:choice_of_parameters}). We note that in this thesis we have considered changes in model parameters such that the change is `slow' relative to the system dynamics. The problem of detecting critical slowing down, and therefore detecting tipping points, in systems with a relatively fast rate of change of critical parameters is a potentially interesting topic for future work, and may be related to the fast-growing field of \emph{rate-induced tipping points}. 
	
	\item {\bf Application of methods to other toy models.} In chapter~\ref{chap:2D} we have used three dynamical systems as examples. It is, however, always possible to apply the techniques presented here to more systems with different types of tipping. Indeed, when meteorological data is presented the type of tipping is often unknown and it is therefore useful to have a variety of techniques available that could be compared to infer the possible nature of the critical transition. In future work the experiments here may be repeated in a variety of different systems and this work may go some way towards the log-term goal of classifying tipping points based on EWS indicator behaviour.
	
	\item {\bf Systems of more than two variables.} The work presented in chapter~\ref{chap:2D} considered only systems of two variables. We consider it would be an interesting project to extend this to systems of a greater number of variables. In particular, the justification of the EOF method in section~\ref{sec:2D:EOF-justification} (page~\pageref{sec:2D:EOF-justification}) used a general two-dimensional system as the basis for the analysis. This analysis could be extended to arbitrary dimensions, although we hypothesise that this will have little effect on the general conclusion.
	
	\item {\bf Estimation of the error in EOF eigenvectors.} In section~\ref{sec:2D:EOF-justification} (page~\pageref{sec:2D:EOF-justification}) where the effectiveness of EOF method is analysed in a tipping-points context, a number of simplifications and generalisations are made. In particular, we did not calculate an error term by considering the variance in the EOF eigenvectors since only the mean value was considered. It does not appear possible to provide this error term without making further simplifications to the equations, but it may be possible to provide an estimate of the error and this will be an important improvement if it results from future work on these equations.
	
	\item {\bf Further investigation of the PS indicator in periodic systems.} In this thesis we have shown that the PS indicator is insensitive to added periodicity in time series when that would cause only a small number of spikes in the power spectrum within the measured frequency range, and possibly none if the frequency of the periodic function lies outside this range. In the application to tropical cyclones (chapter~\ref{chap:cyclones}) the PS indicator, whilst it did not provide a convincing EWS due to the short time series, did appear insensitive to the removal of the 12-hourly oscillations. However, the study of the near-homoclinic system in chapter~\ref{chap:2D} did not provide convincing results. We believe there is scope for further investigation of these effects.
	
	\item {\bf A better choice of application for the PS indicator.} The tropical cyclone example of chapter~\ref{chap:cyclones} is presented as an example of a system with periodicity on similar timescales to the tipping point. This, in theory, provides a good test of the PS indicator properties as mentioned in the previous point. However, as has already been noted, the short time series and small ensembles available made for unconvincing results. A future study involving another periodic time series, but with better available data, would be an interesting project.
	
	\item {\bf A more suitable choice of application for EWS techniques on a 2D field.} We have commented already that the choice of the tropical cyclone problem did not provide a particularly good opportunity for testing EWS techniques which require multiple time series datasets spread over a 2D field (see section~\ref{sec:TC:field_data}, page~\pageref{sec:TC:field_data}). Few locations were identified where multiple weather stations were available in the same geographic region. A coastal area was selected where nine large  hurricanes made landfall in the time period of the available data. This area contained $65$ weather stations with available data, although  only between $21$ and $48$ stations were used in each case depending on the hurricane. In a future study the same methods could be applied to a system for which many hundreds of time series datasets are available at regular, gridded locations over a 2D field — possibly model data or satellite reanalysis data. One interesting subject which was explored preliminarily in this thesis is the dynamics of the \emph{greening of the Sahara}. We have looked into the work of, for example, \cite{bathiany2013p1} which uses model data with high spacial density to detect spacial correlations between geographic areas at discrete points in time. The same model with a higher time resolution is being developed, which would allow the application of our higher-dimensional EWS methods and therefore allow us to compare these methods to the spacial-correlation methods used previously. The data from this model is unavailable at the time of writing but this will make an interesting further study along the same lines as this thesis, possibly giving a better example than the tropical cyclone problem used here. 
\end{enumerate}

%\section{Final remarks}


